\documentclass[12pt,a4paper]{article}

% --- Настройка языка и шрифтов для LuaLaTeX ---
\usepackage{fontspec}
\defaultfontfeatures{Ligatures=TeX,Scale=MatchLowercase}

\setmainfont{Alice-Regular}[
    Path=font/,
    Extension=.ttf,
    BoldFont = Alice-Regular,
    ItalicFont = Alice-Regular,
    BoldItalicFont = Alice-Regular,
    BoldFeatures = {FakeBold=1.6},
    ItalicFeatures = {FakeSlant=0.2},
    BoldItalicFeatures = {FakeBold=1.6,FakeSlant=0.2}
]

\newfontfamily\cyrillicfont{Alice-Regular}[
    Path=font/,
    Extension=.ttf,
    BoldFont = Alice-Regular,
    ItalicFont = Alice-Regular,
    BoldItalicFont = Alice-Regular,
    BoldFeatures = {FakeBold=1.6},
    ItalicFeatures = {FakeSlant=0.2},
    BoldItalicFeatures = {FakeBold=1.6,FakeSlant=0.2}
]

\usepackage{polyglossia}
\setmainlanguage{russian}
\setsansfont{TeX Gyre Heros}
\setmonofont{TeX Gyre Cursor}
\AtBeginDocument{\renewcommand{\textnumero}{{\sffamily №}}}

\usepackage{graphicx}
\usepackage{amsmath, amsfonts, amssymb, amsthm}
\usepackage{mathtools}
\usepackage[protrusion=false]{microtype}
\usepackage{unicode-math}
\usepackage{array}
\usepackage{float}
\usepackage{caption}
\usepackage{enumitem}

\usepackage[paper=letterpaper,
            top=20mm,
            bottom=20mm,
            left=10mm,
            right=10mm,
            footskip=0.4in 
]{geometry}

% --- Колонтитулы ---
\usepackage{fancyhdr}
\pagestyle{fancy}
\fancyhf{}
\setlength{\headheight}{15pt}

\renewcommand{\headrulewidth}{0.4pt}
\renewcommand{\footrulewidth}{0.4pt}

\newcommand{\DocNumber}{1}
\newcommand{\DocTitle}{Аналитическое исследование
механизмов, состояния и
результативности импортозамещения
в производстве вычислительной
техники и периферийных устройств}
\newcommand{\Discipline}{ЭВМ и периферийные устройства}

\title{
	Министерство науки и высшего образования 
    
    Российской Федерации
    
    Национальный исследовательский университет «МИЭТ»
    
    Институт интегральной электроники имени академика 
    
    К. А. Валиева

	\vspace{2em}
}
\author{
}
\date{}

\begin{document}

\maketitle
\thispagestyle{fancy}

\begin{center}
	\Large Эссе \textnumero~$\DocNumber$\\
	по дисциплине «\Discipline»\\
	«\DocTitle»
\end{center}

\vspace{15em}

\begin{flushright}
	Выполнил студент группы ЭН-$35$\\
	Андреев И. А.
\end{flushright}

\cfoot{Москва 2026}

\pagebreak

\fancyhead[L]{\DocTitle}
\fancyhead[R]{\Discipline}
\fancyfoot[C]{\thepage}

\begin{center}
    \Large\textbf{Эссе \textnumero~$\DocNumber$}\\
    \large\textbf{\DocTitle}
\end{center}

\section{Введение}

Политика замещения импорта в сфере электронно-вычислительной 
техники (ЭВМ), микроэлектроники и периферийного оборудования 
перестала рассматриваться исключительно как инструмент догоняющего 
развития развивающихся экономик середины XX века и трансформировалась 
в стратегическую основу национальной безопасности ведущих 
индустриальных государств. Современный этап глобальной 
реконфигурации обусловлен обострением геополитического 
противостояния технологических лидеров, уязвимостью 
трансконтинентальных логистических маршрутов и критической 
зависимостью базовой цифровой инфраструктуры от закрытой 
компонентной базы и проприетарного микрокода ограниченного пула 
транснациональных корпораций.

\subsection{Международный опыт}

Наиболее масштабный и нормативно-оформленный опыт высокотехнологичного
импортозамещения демонстрирует КНР. 
Центральным звеном экономической стратегии Пекина выступает 
государственная программа Xinchuang (инновации в сфере прикладных 
информационных технологий), нацеленная на последовательное 
устранение иностранного аппаратного и программного обеспечения из 
критической информационной инфраструктуры. Программа реализуется в 
рамках трех последовательных фаз: первоочередная замена 
технологического стека в государственных учреждениях и партийных 
структурах, последующий охват системообразующих государственных 
корпораций в энергетике, телекоммуникациях и финансовом секторе, 
и финальное масштабирование на коммерческий и частный сегменты.

Согласно фундаментальным исследованиям института Bruegel 
(Алисия Гарсия-Эрреро, Полин Вейль), полупроводники на протяжении 
длительного времени оставались основной статьей импорта Китая, 
превосходя по финансовым затратам сырую нефть. Санкционное давление 
со стороны США, выразившееся во внесении национальных предприятий в 
списки экспортного контроля, привело к запуску 
Национального инвестиционного фонда интегральных схем, 
совокупный объем вложений которого превысил $150\text{ млрд}$ 
долларов, а третий транш достиг $344\text{ млрд}$ юаней с целевым 
фокусом на производство оборудования литографии, средств 
автоматизации проектирования (EDA) и высокопроизводительной памяти HBM.

Статистические данные Американо-китайской комиссии по экономике 
и безопасности (USCC) и Главного таможенного управления КНР 
подтверждают структурный характер замещения: физический объем 
импорта интегральных схем сократился на $24\%$ в период с $2021$ 
по $2023$ год, а общая стоимость ввозимых полупроводников в $2023$ году 
упала на рекордные $15{,}4\%$. Как подчеркивают аналитики Центра 
стратегических и международных исследований (CSIS), западный 
экспортный контроль непреднамеренно консолидировал внутренних 
китайских потребителей чипов и производителей аппаратуры, преодолев 
прежнее сопротивление корпоративного сектора, предпочитавшего 
зарубежные архитектуры.

Альтернативную модель демонстрирует Индия, ориентирующаяся не на 
технологическую автаркию, а на принудительную локализацию мировых 
сборочных и компонентных цепочек через стимулирование спроса и 
субсидирование предложения. Ключевым драйвером выступает инициатива 
связанных с производством стимулов (Production Linked Incentive, PLI), 
в частности программа PLI 2.0 for IT Hardware с бюджетом 
$17\,000\text{ крор}$ рупий, охватывающая выпуск ноутбуков, серверов, 
планшетов и моноблоков. Индийский механизм предусматривает выплату 
прямой государственной премии в размере $3\text{--}6\%$ на прирост 
чистых продаж локализованной продукции в течение шести лет. Если 
ранее Индия внутренним производством закрывала лишь $10\text{--}20\%$ 
импорта портативных компьютеров, то развертывание схемы PLI совместно 
с субсидированием компонентного производства (Electronics Component 
Manufacturing Scheme, ECMS) позволило привлечь транснациональных 
контрактных сборщиков и довести долю добавленной стоимости внутри 
страны до $23\%$.

\subsection{Специфика российской модели технологического суверенитета}

В отличие от китайской стратегии, опирающейся на гигантский 
масштаб внутреннего рынка и триллионные государственные фонды, и 
индийской модели, интегрирующейся в глобальные цепочки западных 
OEM-производителей, российская политика импортозамещения формировалась 
в условиях жестких экстерриториальных санкций, разрыва логистических 
связей и внезапного ухода глобальных IT-вендоров. Главным целевым 
ориентиром стало обеспечение непрерывности функционирования 
критической информационной инфраструктуры. При этом базовым 
регуляторным инструментом выступили не прямые экспортные премии, а 
жесткий национальный режим государственных закупок, подкрепленный 
балльной системой подтверждения производства в соответствии с 
Постановлением Правительства РФ \textnumero~$719$.

\section{Теоретико-методологические парадигмы и технологические 
переделы производства}

\subsection{Ограничительно-протекционистская и стимулирующая 
инновационная модели}

В фундаментальной макроэкономической теории концепция 
импортозамещающей индустриализации (Import Substitution 
Industrialization, ISI) разделяется на две ведущие парадигмы: 
ограничительно-протекционистскую и стимулирующую инновационную. 
Ограничительная модель концентрируется на создании искусственных 
регуляторных и тарифных барьеров для внешней продукции посредством 
прямых эмбарго, заградительных таможенных пошлин и квотирования 
государственных закупок. Данный подход создает гарантированную 
рыночную нишу для внутренних предприятий, однако при отсутствии 
конкурентной среды и смежных технологических переделов неизбежно 
приводит к инфляционному давлению, деградации стандартов качества 
и технологическому отставанию.

Стимулирующая парадигма базируется на государственной компенсации 
затрат на научно-исследовательские разработки, льготном 
налогообложении центров проектирования, субсидировании процентных 
ставок по кредитам на технологическое перевооружение и формировании 
кластерных экосистем. Как показывают исследования К. Круглова 
(Финансовый университет при Правительстве РФ), устойчивый результат 
в высокотехнологичных секторах достижим исключительно при гармонизации 
обоих подходов, когда административное отсечение импорта строго 
синхронизировано по времени с физическим вводом в эксплуатацию 
замещающих производственных мощностей.

\subsection{Пять базовых технологических переделов вычислительной 
техники}

Специфика вычислительной техники и периферийных систем требует четкой 
дифференциации производственных этапов радиоэлектронной 
промышленности. Процесс локализации сложного цифрового оборудования 
структурируется по пяти базовым технологическим переделам:

Первый передел охватывает крупноузловую отверточную сборку, при 
которой готовые импортированные 
модули монтируются в готовый корпус, что не создает устойчивых 
локальных инженерных компетенций и характеризуется минимальной 
добавленной стоимостью.

Второй передел предполагает мелкоузловую сборку с освоением 
технологии поверхностного монтажа электронных компонентов на 
печатные платы. На этом уровне предприятие самостоятельно 
осуществляет нанесение паяльной пасты, высокоскоростную установку 
поверхностно-монтируемых компонентов, пайку в конвекционных печах, 
оптический и рентгеновский контроль качества.

Третий передел включает производство пассивной и конструкционной 
компонентной базы, в том числе изготовление многослойных 
несмонтированных печатных плат высокой плотности межсоединений, 
литье пластиковых элементов корпусов под давлением, изготовление 
алюминиевых радиаторов охлаждения и кабельных сборок.

Четвертый передел заключается в корпусировании, финальной сборке и 
тестировании интегральных схем, включая монтаж кремниевых кристаллов 
на подложку методами ультразвуковой микросварки проводников или 
перевернутого кристалла, компаундирование и валидацию логики микросхем.

Пятый передел представляет собой фундамент полупроводниковой 
индустрии --- фабричное производство кремниевых полупроводниковых 
пластин, включающее операции эпитаксиального наращивания, глубокого 
фотолитографического экспонирования, ионной имплантации и 
химико-механической планаризации.

\subsection{Механизм управления: балльная система Постановления 
\textnumero~719}

Механизмом управления этой многоуровневой структурой со стороны 
государства выступает балльная система подтверждения производства 
промышленной продукции на территории страны, закрепленная 
постановлением Правительства РФ \textnumero~719. Вместо бинарной 
оценки («отечественный» или «зарубежный» товар) регулятор присваивает 
изделию фиксированное количество баллов за каждую фактически 
освоенную технологическую операцию: разработку конструкторской 
документации, реализацию микрокода BIOS, пайку материнской платы, 
монтаж периферийных разъемов, использование локального блока питания 
или интеграцию отечественного микроконтроллера. Ежегодное повышение 
минимального проходного балла формирует долгосрочный нормативный 
стимул, побуждающий частные предприятия непрерывно реинвестировать 
прибыль в углубление производственных цепочек.

\section{Анализ текущего состояния отрасли в Российской Федерации}

Современное состояние сегмента вычислительной техники и периферийных 
элементов в Российской Федерации демонстрирует выраженную внутреннюю 
асимметрию: динамичное развитие контрактных сборочных площадок и 
поверхностного монтажа плат сочетается с критической зависимостью от 
импорта активной электронной компонентной базы и функциональных узлов.

В сегменте серверного оборудования и высокопроизводительных систем 
зафиксированы существенные сдвиги в структуре поставок для 
государственного сектора. Согласно отраслевому мониторингу TAdviser, 
доля серверов отечественных брендов в закупках государственных 
органов к $2024$ году достигла $22\%$. В сегменте систем хранения 
данных, подпадающем под действие нормативных актов об обеспечении 
технологической безопасности критической информационной инфраструктуры, 
доля оборудования российского происхождения в общей структуре 
находящихся в эксплуатации мощностей составила $13{,}7\%$.

Однако разрыв производственных связей с мировыми контрактными 
полупроводниковыми фабриками кардинально усложнил архитектурную 
независимость. Прекращение тайваньской фабрикой TSMC в $2022$ году 
контрактного выпуска микропроцессоров линеек «Эльбрус» 
(АО «МЦСТ», проектные нормы $16$--$28$ нм) и «Байкал» 
(АО «Байкал Электроникс», нормы $16$--$28$ нм) вызвало острый 
дефицит отечественных чипов на открытом рынке. Если в $2021$ году 
государственные ведомства заявляли консолидированную двухлетнюю 
потребность в $1{,}7\text{ млн}$ отечественных процессоров, то 
физическое отсутствие складских остатков сделало эти планы 
нереализуемыми.

В ответ на дефицит процессоров Правительство РФ ввело временную 
двухуровневую систему реестра промышленной продукции, разрешив 
присваивать статус отечественных серверам и персональным компьютерам, 
построенным на зарубежных процессорах (включая архитектуры Intel и 
AMD), при условии выполнения всех остальных критериев локализации 
материнских плат, схемотехники и встроенного программного обеспечения, 
сохраняя при этом приоритет в баллах за конфигурациями на 
отечественных кристаллах.

\subsection{Сегмент периферийного оборудования и печатающей техники}

В сегменте периферийного оборудования, традиционно характеризовавшемся 
тотальным импортом, наблюдается перестройка локального выпуска 
мониторов, автоматизированных рабочих мест и печатной техники. 
Группа компаний «Бештау» развернула промышленные комплексы в 
Ростове-на-Дону и Ессентуках с плановыми мощностями до 
$350\text{ тысяч}$ мониторов и сотен тысяч клавиатур и компьютерных 
мышей в год. Предприятие освоило собственное изготовление печатных 
плат скалеров, литье корпусных деталей и пайку разъемов. Вместе с 
тем жидкокристаллические и OLED матрицы, составляющие до 
$60\text{--}70\%$ себестоимости дисплеев, полностью закупаются в 
странах Азии в силу отсутствия в РФ рентабельных фабрик по производству 
дисплейных панелей.

\subsection{Корректировка государственных нормативов цифровой трансформации}

Осознание производственных ограничений микроэлектроники вынудило 
регулятора скорректировать плановые индикаторы цифровой трансформации. 
Исходные целевые ориентиры предполагали достижение $80\%$ доли 
отечественной вычислительной техники в государственных структурах к 
$2030$ году. Стратегические направления, принятые Правительством, 
установили порог присутствия отечественного оборудования в 
государственном секторе к 2030 году на уровне $58{,}7\%$, фактически 
признавая необходимость сохранения до $41{,}3\%$ импортного парка 
аппаратуры для поддержания стабильности инфраструктуры.

Доля отечественных серверов в закупках госсектора увеличилась с менее 
чем $6\%$ в $2020$ году до $22{,}0\%$ по итогам $2024$ года, а доля 
отечественных систем хранения данных в активной эксплуатации выросла с уровня менее 
$5\%$ до $13{,}7\%$. При этом общая доля отечественной 
радиоэлектроники в закупках госсектора превысила $60{,}0\%$ к 
$2025$ году при скорректированном плановом нормативе в $58{,}7\%$, а 
в сегменте электронных компонентов доля отечественной продукции 
составляет порядка $26{,}7\%$ (с целевым показателем до $50\%$ в 
отдельных вычислительных узлах).

\section{Комплекс действий акторов по организации импортозамещения}

\subsection{Нормативно-регуляторные инициативы государства}

Государственные институты сосредоточили усилия на устранении 
нормативных лазеек, позволявших заказчикам обходить требования 
национального режима государственных закупок (ФЗ-$44$ и ФЗ-$223$). 
Ранее распространенной практикой являлось формирование смешанных 
тендерных лотов, объединявших в одной закупке вычислительное 
оборудование, присутствующее в реестре Минпромторга, с позициями, 
не имеющими отечественных аналогов. В подобных случаях нормативный 
запрет на приобретение иностранной техники переставал действовать 
на всю номенклатуру закупки, что позволяло заказчикам полностью 
приобретать импортные комплексы.

Принятые изменения в постановление Правительства \textnumero~$1875$ 
нормативно запрещают смешивание позиций с различными режимами допуска, 
разделяя закупки на жесткие перечни прямого эмбарго и режим «третий 
лишний» с $15\%$ преференцией по цене для отечественной 
продукции.

\subsection{Инвестиционные и производственные стратегии частного бизнеса}

Отечественный частный бизнес переориентировал 
инвестиционные потоки на расширение собственных контрактных 
площадок. Промышленные компании ввели в строй автоматизированные 
линии поверхностного монтажа плат, системы селективной пайки и 
роботизированные сборочные конвейеры.

Существенный объем инвестиций частных компаний направлен на создание 
доверенного низкоуровневого программного обеспечения: разработку 
суверенных микрокодов BIOS (UEFI), прошивок для сервисных 
микроконтроллеров управления платой и архитектуры контроллеров 
дисковых массивов. 

Чтобы минимизировать рыночные риски, частные производители 
реализуют стратегию гибкой адаптации, сохраняя параллельные 
аппаратные линейки: сертифицированные защищенные платформы на 
базе требований Указа Президента РФ \textnumero~$166$ для объектов 
критической информационной инфраструктуры и коммерческие серверные 
конфигурации на доступных азиатских или параллельно импортируемых 
чипах для открытого корпоративного сектора.

\subsection{Роль академического сообщества и инженерных коллективов}

На уровне научно-исследовательских коллективов, инженеров и сообществ 
разработчиков решаются задачи локализации проектной документации и 
адаптации открытых технологий. Университетские лаборатории и проектные 
дизайн-центры ведут комплексные работы по реверс-инжинирингу 
зарубежной микроэлектроники и периферийных контроллеров, включающие 
декапсуляцию микросхем, послойный электронно-микроскопический анализ 
топологии кристаллов, восстановление электрических схем и воссоздание 
микропрограммных алгоритмов.

Параллельно инженерные группы наращивают участие в разработке ядер 
на открытой архитектуре команд RISC-V, стремясь сформировать 
независимый процессорный задел для контроллеров накопителей, 
коммутаторов и периферийных устройств в обход лицензионных 
ограничений экосистем ARM и x86. В программном сегменте инженерные 
коллективы обеспечивают интеграцию аппаратных платформ отечественных 
ПК и серверов с дистрибутивами российского системного ПО 
(Astra Linux, РЕД ОС, Базальт СПО), создавая локализованные 
драйверы печати, сетевых адаптеров и специализированных плат 
ввода-вывода.

\section{Оценка результативности}

Анализ практики реализации программ импортозамещения позволяет 
выделить устойчивые положительные результаты наряду со скрытыми 
издержками, технологическими искажениями и рисками квазизамещения.

\subsection{Полезные эффекты и структурные достижения}

К категории полезных и устойчивых результатов относится формирование 
развитой национальной инфраструктуры поверхностного монтажа печатных 
плат и финальной сборки радиоэлектронной продукции. Созданные частным 
бизнесом автоматизированные производственные мощности гарантируют 
физическую сборку серверов, рабочих станций, мониторов и контроллеров 
внутри страны. 

Фундаментальным достижением стало устранение критических уязвимостей 
на микропрограммном уровне вычислительных систем: внедрение 
проверенных отечественных реализаций микрокодов BIOS и сервисных 
процессоров BMC предотвратило возможность скрытого деструктивного 
воздействия через встроенные механизмы аппаратного управления 
мировых вендоров. 

В периферийном сегменте полезным результатом 
явилась локализация базовых узлов питания, плат управления и 
микропрограммного обеспечения печатных и мониторных систем, что 
закрыло потребности государственных ведомств в типовом оборудовании.

\subsection{Скрытые издержки, системные аномалии и технологический долг}

Масштабное исследование 
аналитического центра The Russia Program фиксирует, что стремление к 
суверенитету не устранило внешнюю уязвимость цифрового сектора, а 
сместило ее вектор в сторону зависимости от КНР. Анализ практики 
корпоративного перехода, проведенный компанией K$2$ Cyber Security, 
показал, что для $30\%$ предприятий импортозамещение выразилось 
исключительно в замене западных компонентов и программных продуктов 
на китайские решения. В результате отечественная вычислительная 
техника сохраняет критическую зависимость от зарубежных поставок 
дискретной логики, пассивных компонентов, контроллеров и памяти.

Слабая проработка контрольных процедур на ранних этапах 
импортозамещения привела к широкому распространению практики 
крупноузловой сборки готовых азиатских платформ с формальной 
русификацией интерфейсов и сменой логотипа на корпусе. Подобный 
подход отвлекал государственные ресурсы, создавая фиктивные 
реестровые решения без реального трансфера конструкторских компетенций.

Малосерийный характер российского производства аппаратных комплексов 
исключает реализацию эффекта экономии на масштабе, присущего фабрикам 
Юго-Восточной Азии. Себестоимость локализованных систем хранения 
данных и серверов оказывается на $20\text{--}50\%$ выше зарубежных 
оригиналов при сопутствующих задержках логистики параллельного 
импорта комплектующих до $8\text{--}12$ недель. Вынужденный перевод 
информационных систем на менее производительную и дорогую аппаратную 
базу приводит к снижению общей операционной эффективности потребителей 
и формированию отложенного технологического долга.

Разработка собственных процессорных архитектур (fabless-модель) без 
наличия доступных суверенных кремниевых фабрик необходимого 
топологического уровня оказалась тупиковой стратегией в условиях 
жестких санкций. Попытки перенести производство отечественных чипов 
линеек «Байкал» и «Эльбрус» на китайские контрактные фабрики SMIC и 
Hua Hong столкнулись с фактическим отказом последних из-за высоких 
рисков подпадания под вторичные санкции со стороны регуляторов США.

\section{Выводы}

Комплексный анализ мирового опыта и текущего состояния 
радиоэлектронной промышленности Российской Федерации позволяет 
сформировать систему выводов и прикладных рекомендаций по оптимизации 
государственной политики в секторе производства вычислительной техники 
и периферийных элементов.

Попытка форсированного построения замкнутой технологической автаркии 
полного цикла в области производства сложной вычислительной техники 
экономически нецелесообразна ввиду колоссальной капиталоемкости 
микроэлектронных фабрик и ограниченной емкости внутреннего рынка 
сбыта. Корректировка целевого показателя доли отечественной 
вычислительной техники в государственных структурах до $58{,}7\%$ к 
$2030$ году свидетельствует о прагматизации регуляторных подходов и 
признании объективной необходимости сохранения технологических мостов 
с дружественными индустриальными юрисдикциями.

В условиях долгосрочной блокировки доступа отечественных 
дизайн-центров к передовым контрактным зарубежным фаундри (TSMC) 
на базе архитектур x$86$ и проприетарного ARM, единственным 
устойчивым направлением развития вычислительной техники выступает 
концентрация государственных субсидий и частных инвестиций вокруг 
международной открытой архитектуры RISC-V. Опыт Китая демонстрирует, 
что разработка микроконтроллеров и процессоров на базе открытых 
стандартов позволяет формировать независимый микропрограммный стек 
без риска внезапного отзыва коммерческих лицензий держателями 
западной интеллектуальной собственности.

Для исключения дрейфа зависимости от одних внешних монополистов к 
другим государственные механизмы подтверждения происхождения продукции 
(постановление Правительства \textnumero~$719$) должны непрерывно 
повышать нормативные требования к локализации физической компонентной 
базы. Опыт индийской схемы ECMS подтверждает эффективность 
стимулирования спроса именно на отечественную электронно-компонентную 
базу: преференции на государственных торгах должны предоставляться 
не за факт монтажа печатной платы иностранными чипами, а 
пропорционально доле установленных отечественных полупроводниковых и 
пассивных компонентов. Такой подход позволит трансформировать 
поверхностную контрактную сборку в устойчивую и технологически 
суверенную отрасль аппаратного приборостроения.

\end{document}
